\chapter{Foundations and Related Work}\label{foundations-and-related-work}

The closest cross-artifact verification work was identified through a targeted
Google Scholar search on 30 July 2026. The search combined terms for
natural-language software artifacts, GUI or screenshot artifacts, and
multimodal verification. It was exploratory rather than a systematic
literature review and was supplemented by references from the identified
studies. The comparison focuses on work that connects textual
software-engineering artifacts with visual UI artifacts through a verification,
consistency, or grounding decision. The following sections also draw on the
supporting literatures on requirements quality, traceability, GUI testing,
interaction trajectories, visual grounding, and abstention.

\section{Requirements as Verification Contracts}\label{requirements-as-verification-contracts}

A software requirement states an obligation that an implementation is expected
to satisfy. Requirements may describe functionality, quality properties,
constraints, interactions, or externally visible outcomes. In practice, their
role is broader than that of natural-language documentation. They coordinate
expectations between stakeholders and provide the reference against which
design, implementation, and testing decisions are justified. Verification
depends on the implemented system as well as the precision and observability of
the requirement.

Natural-language requirements are attractive because stakeholders can read and
write them without committing to a formal notation. The requirements studied
here are natural-language artifacts and therefore retain ambiguities arising
from wording, domain context, and interpretation. Linguistic sources include
vague terms, implicit references, and underspecified quantifiers (Berry,
Kamsties, and Krieger, 2003; Gervasi et al., 2019). Inspecting an implementation
cannot by itself resolve an underspecified target. If a requirement asks for
``all relevant results'' without defining relevance, even complete access to the
UI would not establish a unique expected outcome.

Requirements retain their stated scope throughout the evaluation. Quantifiers
such as ``all,'' ``only,'' ``always,'' and ``closest'' can be supported when the
interface presents a bounded set as complete. Open or ambiguous domains require
uncertainty rather than an assumption that the observed cases are
representative. The exact decision rules are defined in Section 3.4.

UI requirements occupy a subset of software requirements.
They may constrain visible content, available controls, navigation, validation
feedback, state changes, or the presentation of results. Some are directly
observable in a single state, while others become observable only through a
sequence. A requirement that a selected item appears in a cart needs at least a
selection state and a later cart state. A requirement that a menu is reachable
without authentication requires the navigation path and the absence of an
intervening sign-in wall. This temporal structure motivates the use of ordered
screenshot flows rather than isolated images.

\section{Requirements Traceability}\label{requirements-traceability}

Traceability links requirements to the artifacts that refine, implement, or
verify them. \textcite{clelandhuang2014} characterize software traceability as a
persistent challenge because projects contain heterogeneous artifacts, links
change as systems evolve, and manual maintenance is expensive. Classical links
may connect requirements to design elements, source code, test cases, defects,
or other requirements. The present task adds a visual trace target: a
requirement decision should be connected to the recorded UI states that support
or contradict it.

This link is not equivalent to a conventional document-retrieval result. A
screenshot can contain the same words as a requirement without demonstrating
the required behavior. Conversely, an icon, selected state, or layout change may
provide relevant evidence without repeating the requirement vocabulary. A
complete trace may also span several steps. The system must therefore preserve
both semantic relevance and temporal position.

Trace quality and classification accuracy are measured independently. Without
a trace, a reviewer has to repeat the complete inspection. A cited screenshot
step and region expose the observation on which the model relied. Correct
labels may still have incorrect traces, and correct traces may still be
interpreted incorrectly.

\section{GUI Testing and Screenshot-Based Verification}\label{gui-testing-and-screenshot-based-verification}

GUI test automation normally executes interactions and checks explicit
assertions against a running system. It can provide strong evidence when the
environment, selectors, input data, and expected states are controlled. In
real-world systems, however, interfaces change, element locators become brittle,
external services introduce nondeterminism, and visually equivalent states may
have different internal representations. \textcite{nass2021}
identify recurring GUI test-automation challenges that have persisted for many
years and argue that some are likely to remain.

The evaluation begins after a flow has been recorded and complements executable
GUI tests. The original application can be unavailable, and no interaction
script has to be replayed. A shared screenshot representation also permits
comparison across heterogeneous websites. The verifier does not execute the
application and therefore cannot observe backend state, network effects, or
paths outside the captured execution. For the Mind2Web benchmark, lexical
screenshot selection additionally uses visible-text candidates extracted from
recorded HTML metadata. This metadata ranks screenshots but is not treated as
fulfillment evidence: the multimodal verdict is based on the selected images
and image-derived OCR hints. The same selector can operate on OCR-only
representations, although screenshot-only inputs may provide a weaker retrieval
signal.

These omissions delimit the claims available from screenshots. A confirmation
message is evidence for the message, not for email delivery. A selected
checkbox says nothing about a later session, and one result page cannot
establish global completeness. The evaluation concerns conformance to the
recorded visual evidence.

\textcite{kretzer2025} provide a close requirements-oriented comparison by
connecting user stories to GUI prototypes and studying whether a story is
represented in the prototype. The present work differs by using ordered
observations of an executed interface, an explicit four-label contract, and
evidence links that may span multiple states.

\textcite{kolthoff2025guispector} move closer to executable requirement verification.
GUISpector actively explores a GUI, records the resulting trajectory, and
returns evidence-backed met, partially met, or unmet decisions for requirements
and their acceptance criteria. This thesis instead holds the observed
trajectory fixed so that evidence selection and decisions under missing states
can be evaluated separately from the quality of an exploration policy. Its
four-label contract also represents insufficient evidence explicitly through
\texttt{ABSTAIN}.

\textcite{massenon2026} study cross-modal verification of mobile-app
bug fixes and demonstrate the broader relevance of combining textual
engineering artifacts with visual evidence. Their task concerns bug-fix
consistency rather than requirement fulfillment, but it motivates treating
multimodal verification as a software engineering problem rather than generic
visual question answering.

\section{Multimodal UI Agents and Interaction Trajectories}\label{multimodal-ui-agents-and-interaction-trajectories}

Multimodal UI agents receive visual and textual observations and select actions
in an interface. Their progress is relevant because both UI action selection and
UI verification require models to connect language with visible state. The
tasks nevertheless have different objectives. An agent can complete a task
without explaining which evidence establishes a requirement, and a verifier
does not need to choose or execute the next interaction.

Mind2Web contains open-ended web tasks, action sequences, and observations from
real websites \autocite{deng2023}. Its original purpose is to evaluate
generalist web agents across tasks, websites, and domains. This thesis reuses a
small set of its ordered trajectories as visual observations, not as an
action-prediction benchmark. The original task description and action sequence
provide context for reconstructing a flow, while the thesis benchmark adds
separately reviewed UI requirements, fulfillment labels, and evidence
annotations.

\section{Visual Grounding and Evidence Localization}\label{visual-grounding-and-evidence-localization}

Visual grounding connects a linguistic expression to a spatial region. In UI
settings, the target may be a control, text block, icon, status indicator, or
compound arrangement. SeeClick trains and evaluates grounding capabilities
specialized for graphical user interfaces and shows the importance of
interface-specific localization \autocite{cheng2024}. UGround similarly treats
GUI grounding as a dedicated capability rather than assuming it emerges
reliably from a general-purpose vision-language model \autocite{gou2025}.

Set-of-Mark prompting overlays numbered candidate regions and asks a
multimodal model to refer to the marks instead of emitting unconstrained pixel
coordinates \autocite{yang2023}. The approach can simplify the mapping from
language to pixels, but its upper bound depends on the proposal generator: a
mark-selection step cannot return a relevant region that was omitted from the
candidate set. Mark density and placement may also obscure the interface.
SeeAct reports that Set-of-Mark was not its strongest strategy for web-agent
grounding and benefited from combining visual information with HTML-derived
candidates \autocite{zheng2024}.

The auxiliary candidate-mark implementation uses OmniParser-derived UI
proposals together with OCR candidates (Lu et al., 2024). The main
label-evaluation matrix instead uses direct model-generated regions and does
not depend on candidate-mark grounding. The auxiliary variant is evaluated as
evidence localization. Coordinate validity covers only geometry; semantic
correctness additionally requires a region that is relevant to the claim and
sufficiently specific for review.

\section{Abstention and Decision-Making under Incomplete Evidence}\label{abstention-and-decision-making-under-incomplete-evidence}

Classification with a reject option permits a system to withhold a concrete
prediction when the expected cost of an error is too high or the available
information is insufficient. \textcite{hendrickx2024} survey reject-option
methods across machine learning, while \textcite{wen2025} discuss abstention in
large language models. The common intuition is applicable here, but the thesis
does not attempt to learn a calibrated confidence threshold.

\texttt{ABSTAIN} is instead a semantic label in the supplied verification contract. It
means that the screenshot flow does not justify a reliable positive or negative
decision. This differs from a missing API output, parsing failure, or skipped
item. Those are coverage failures and must be reported separately. \texttt{ABSTAIN}
also differs from \texttt{NOT\_FULFILLED}. The latter requires visible counter-evidence,
whereas abstention may follow from an omitted result state, hidden property, or
unresolved textual or evidence ambiguity.

False-positive \texttt{FULFILLED} verdicts can end a review prematurely. The evaluation
therefore reports false fulfillment alongside accuracy. An all-abstain system
would avoid many positive errors but provide no useful decisions. Coverage,
per-class performance, and recorded reasons therefore contextualize the
abstention rate.

\section{Synthesis of the Research Gap}\label{synthesis-of-the-research-gap}

Taken together, these research strands establish the relevance of requirements ambiguity, traceability,
GUI automation, multimodal software verification, UI-agent trajectories,
visual grounding, and abstention. The work identified in the targeted search
does not directly evaluate the following contract: given a textual requirement and an already recorded ordered
UI flow, assess its fulfillment, identify the supporting or contradicting
screenshot steps, optionally localize the evidence within those screenshots,
and preserve uncertainty about hidden or missing states.

The reviewed multimodal systems can process screenshots, but the identified work does
not provide the full verification contract studied here. The missing elements
are a controlled formulation that separates requirement semantics, flow
coverage, evidence traces, and label aggregation. Chapter 3 defines this
formulation.

Table \ref{tab:research-gap-comparison} makes the boundary to adjacent research
areas explicit. The central distinction is the combination of an ordered
evidence trace with a fulfillment decision. Explicit abstention and uncertainty
reasons add diagnostic context but are not treated as a separate research
contribution.

\begin{table}[htbp]
  \centering
  \footnotesize
  \setlength{\tabcolsep}{4pt}
  \begin{tabular}{@{}>{\raggedright\arraybackslash}p{0.20\linewidth}>{\raggedright\arraybackslash}p{0.18\linewidth}>{\raggedright\arraybackslash}p{0.18\linewidth}>{\raggedright\arraybackslash}p{0.14\linewidth}>{\raggedright\arraybackslash}p{0.23\linewidth}@{}}
    \toprule
    \textbf{Research area} & \textbf{Primary input} & \textbf{Typical output} &
    \textbf{Ordered evidence} & \textbf{Explicit abstention or uncertainty} \\
    \midrule
    \textbf{Requirements traceability} & Requirements and engineering artifacts & Trace links & Sometimes & Not usually a fulfillment label \\
    \textbf{GUI test automation} & Executable UI and test oracle & Pass/fail and execution trace & Yes & Usually encoded as test outcome \\
    \textbf{Mind2Web-style UI agents} & Task instruction and web trajectory & Next action or task success & Yes & Not the central output \\
    \textbf{Visual UI grounding} & Instruction and UI image & Element or region & Usually no & Localization failure rather than requirement abstention \\
    \textbf{User-story/prototype support} & User story and prototype & Story--prototype consistency & Usually static & Task-specific \\
    \midrule
    \rowcolor{TUMLightGray!30}
    \textbf{This thesis} & \textbf{Requirement and recorded screenshot flow} &
    \textbf{Four-way label plus evidence} & \textbf{Yes} &
    \textbf{Explicit \texttt{ABSTAIN} and uncertainty reasons} \\
    \bottomrule
  \end{tabular}
  \caption{Comparison of adjacent research areas with the verification contract studied in this thesis.}
  \label{tab:research-gap-comparison}
\end{table}
